% ---------------------------------------------------------------------------
% METHODS — section master file.
% \input by main/Manuscript.tex. Each \subsection lives in its own file in this
% folder; this file holds only the section header, the opening paragraph, and
% the running order. Add new subsections as files and \input them below.
%
% ORDER follows the earmold paper (Friedrich & Schoenwiesner 2025): behavioral
% methods first, acoustic methods after Statistical analysis.
%
% WARNING (2026-08-26): main/Manuscript.tex \inputs this file AND still carries
% a full inline copy of the old earmold Methods further down. Until that inline
% block is removed, the section renders TWICE and \label{sec:2} is defined
% twice. See the note at the bottom of this file.
% ---------------------------------------------------------------------------

\section{\label{sec:2} Methods}

Participants had their individual HRTF recorded and trained to localize sounds
with one ear only in a virtual auditory environment (VAE) for 4 consecutive
days. To capture the trajectories of participants' adaptation to the HRTFs
presented in the VAE, participants completed daily localization tests. During
training and testing, HRTFs were presented to one ear only.

% --- behavioral methods ----------------------------------------------------

\subsection{Participants}
% TODO — n, age range, ethics approval, compensation. Recruitment ran on
% 18-40 years / normal hearing / 4 consecutive days. Planned n = 8, raised to
% 16 if the statistics require it.

\subsection{Experimental procedure}
% TODO — the 4-day timeline. Day 1: native binaural reference, donor build,
% baseline A and baseline D. Adaptation days: PRE test, training, POST test.
% Final day: counterbalanced 2x2 (Ear x Side), Williams Latin square.

% ---------------------------------------------------------------------------
% DRAFT SEGMENT — experimental setup
% Standalone; not \input anywhere yet. Stitch into Methods later.
% Code: experiment/localization/Localization_AR.py   (rendering + response)
%       hrtf/binsim/hrtf2binsim.py                   (settings, reverb, hp filter)
%       hrtf/binsim/hrir2mat.py                      (level normalization)
% Room / array / response apparatus are cited to the earmold paper rather than
% re-described — they are unchanged, and every reported recording came through
% record_dome on the existing speaker table (params.txt, all subjects).
% ---------------------------------------------------------------------------

\subsection{Experimental setup}

Recordings, localization tests and training were carried out in the
hemi-anechoic room and 45-loudspeaker array described previously
% \citep{friedrich_adaptation_2025}
, using the same head-mounted laser and electromagnetic motion sensor
(MetaMotionRL, Mbientlab, San Francisco, USA) for visual feedback of head
orientation and for the response. Acoustic measurements used the vertical
midline column of seven loudspeakers, spanning $\pm37.5^\circ$ elevation at
$12.5^\circ$ spacing. Signals were generated and acquired with TDT System 3
hardware (Tucker Davis Technologies, Alachua, USA) at a sampling rate of
48.8\,kHz.

Virtual sound sources were rendered over headphones (DT 990 Pro,
beyerdynamic, Heilbronn, Germany) by real-time convolution in pyBinSim
% \citep{neidhardt_pybinsim_2017}
, which selected the impulse response pair for the listener's current head
orientation and crossfaded between successive filters. Convolution ran in
blocks of 256 samples on a 512-tap direct-sound filter, followed in series by a
late reverberation tail at a direct-to-reverberant ratio of 20\,dB and the
participant's own headphone equalization filter. The gain of the rendered chain is
not fixed by the measurement, so each database was scaled by a single
broadband factor to a common value; presentation
level was thereby made identical across participants, conditions and days,
while interaural and between-direction level differences, being ratios within
a database, were unaffected. Overall level was matched by ear to the
free-field array and the system volume held fixed thereafter.

% --- open items ------------------------------------------------------------
% - SELF-CITATION verified: Friedrich, P., and Schoenwiesner, M. (2025).
%   "Adaptation rate and persistence across multiple sets of spectral cues for
%   sound localization," J. Acoust. Soc. Am. 157(3), 1543-1553,
%   doi:10.1121/10.0036056. Key friedrich_adaptation_2025; still needs adding
%   to the .bib, which does not yet exist for this manuscript.
% - pyBinSim citation not verified. Neidhardt et al., pyBinSim, is the usual
%   one; check what the pfriedrich-hub/pybinsim_tuil fork should be cited as,
%   and whether the fork's deviations need a sentence.
% - The late reverberation tail: where does it come from? Generic or measured?
%   One clause is owed, a reviewer will ask whether it is individual.
% - End-to-end latency of the rendering chain is NOT stated. One block is
%   5.24 ms at 48.8 kHz, but total motion-to-sound latency includes the sensor
%   rate and buffering and has not been measured. Either measure it and give a
%   number, or say nothing -- do not quote the block period as the latency.
% - Headphone model: DT 990 Pro confirmed for all reported participants
%   (MYSPHERE appears in the code as an alternative, unused here).
% - LEVEL JUSTIFICATION REWORDED 2026-08-26. It used to read "impulse responses
%   derived from directional transfer functions carry no absolute scale", which
%   rested on the DTF framing. A properly referenced HRTF DOES have a defined
%   magnitude (relative to free field), so that justification no longer holds.
%   The 6.5-15.2 dB chain-gain spread across databases is therefore NOT
%   explained by "no absolute scale" -- it must come from the processing
%   (windowing, onset alignment, low-frequency extrapolation) or from a real
%   between-subject difference, and 8.7 dB is too large to be the latter.
%   WORTH CHECKING: if the HRTFs are properly referenced, where does the spread
%   come from? The current wording is deliberately neutral and does not claim a
%   cause. Related: project_binsim_level_normalization.
% - Room dimensions deliberately not repeated; if a reviewer wants the paper
%   standalone, restore "6.4 x 4.4 x 2.3 m" to the first sentence.


\subsection{Localization task}
% TODO — sector grid (7 deg az, 14 deg el), el +-35, 3 targets per sector,
% min separation 20 deg, midline excluded from one-sided blocks; 75 trials per
% one-sided block, 150 for the binaural full-field reference. Also the
% free-field control: same task as Friedrich & Schoenwiesner (2025) but
% restricted to the seven midline loudspeakers, 3 targets each (21 trials),
% because the midline is the only direction at which the virtual rendering
% reproduces the measurement rather than the spherical-head model.

% ---------------------------------------------------------------------------
% DRAFT SEGMENT — localization test stimuli
% Standalone; not \input anywhere yet. Stitch into Methods later.
% Backing analysis + parameter sweeps: docs/stimulus_spectral_variation.md
% Code: experiment/localization/localization_helpers/stimulus.py
% ---------------------------------------------------------------------------

\subsection{Localization stimuli}

Localization stimuli were 225\,ms trains of five 25\,ms pink-noise bursts
separated by 25\,ms silent gaps (10\,ms raised-cosine onset and offset ramps,
5\,ms at the gap edges), cut from a single continuous noise realization. In the
fixed-spectrum condition the train was used unmodified, so its expected
spectrum was the same on every trial and only the realization noise varied
(SD of the 1/6-octave-smoothed log spectrum: 0.4\,dB across trains, 0.8\,dB
across the bursts within a train).

In the variable-spectrum condition each train was recolored, before gapping, by
a random spectral envelope drawn anew on every trial and shared by all five
bursts. Natural sound sources differ from one another mainly in broad spectral
coloration, which the auditory system must discount in order to read the finer,
elevation-dependent pattern imposed by the pinna; the envelope was built to
impose exactly that kind of variation. It was a smooth random function of log
frequency (0.5--16\,kHz), band-limited to ripple densities below 0.5 ripples
per octave with an rms of 3\,dB (median peak-to-trough depth 10\,dB), applied as
a zero-phase gain. The cutoff follows Macpherson and Middlebrooks
% \citet{macpherson_vertical-plane_2003}
, who found that localization was disrupted by spectral ripple at densities of
0.5--2 ripples per octave but not by coarser variation, and concluded that
broad-scale spectral structure is discounted despite contributing most to the
shape of directional transfer functions. Each trial was
therefore strongly recolored while the 0.5--2 ripples/octave band, in which the
pinna notch shifts with elevation, was left largely untouched. The manipulation
thus tests whether what was learned during training is tied to the spectrum of
the training stimulus or generalizes across source spectra. It does not
separate a perceptual remapping from an explicitly learned association: a
listener who had learned to read the preserved band would succeed under either
account. Timing, bandwidth and level were unchanged, and the envelope of every
trial was stored.

Across 400 realizations the variable stimulus varied with an SD of 2.4\,dB,
against 0.4\,dB for the fixed-spectrum stimulus. Within the 0.5--2
ripples/octave band the elevation-dependent variation of the participants'
head-related transfer functions (SD across elevation at fixed azimuth) exceeded
the trial-to-trial variation of the source spectrum by a factor of 7.5; below
0.5 ripples/octave the two were comparable (factor 1.0).

% NB: every number in this file scales linearly with the envelope rms, so
% changing the depth is a proportional edit throughout. Values at the earlier
% 8 dB setting, for reference: 27 dB median peak-to-trough, 6.4 dB stimulus SD,
% 2.8 in the cue band, 0.4 below it. Settle this against the self-test
% (protocols/dev/stimulus_check.py cells 7-9) before submission.

\subsubsection*{Choice of envelope depth}

Envelope depth is bounded from both sides, and the two bounds are what fixed
it. It must be large enough that the coarse-scale coloration of the source is
no longer a reliable guide to elevation, and small enough that the elevation
cue itself survives. The upper bound is the more concrete of the two.
Macpherson and Middlebrooks
% \citet{macpherson_vertical-plane_2003}
report degradation of vertical-plane localization at 1 ripple per octave only
for peak-to-trough depths at or above 20\,dB, and at coarser densities even
40\,dB left performance largely intact; we therefore treated 20\,dB
peak-to-trough as a ceiling and stayed well below it. Two further
considerations argue for the same direction. The band limit is not sharp---a
function band-limited on a finite log-frequency axis has skirts---so a deeper
envelope leaks into the 0.5--2 ripples/octave region even though no energy is
assigned there, and the depth of the elevation cue itself varies between ears
and listeners, so a setting chosen against one head-related transfer function
must leave headroom for shallower ones.

The lower bound is set by the purpose of the manipulation. In the fixed-spectrum
condition the sound at the eardrum carries a near one-to-one map from absolute
spectrum to elevation, so a listener can in principle solve the task by
matching the overall spectrum to a stored template, without separating source
from filter. The envelope must make that route unreliable, which requires its
trial-to-trial variation to be at least comparable to the elevation-dependent
variation of the head-related transfer function at the same spectral scales. At
3\,dB rms the two are of the same magnitude below 0.5 ripples per octave
(factor 1.0), so overall coloration no longer specifies elevation, while the
cue band retains a 7.5:1 margin. Deeper envelopes would make the source
dominate rather than merely match at coarse scales, but only at the cost of
the margin in the cue band, and of exceeding the depth at which disruption has
been reported.

Variation was confined to the coarse band rather than spread across all
spectral scales for the same reason: variation inside the cue band would
contaminate the very structure being read, and contributes nothing to defeating
a template-matching strategy that the coarse band does not already achieve. It
was retained as an adjustable parameter and set to zero throughout, so that
deliberately contaminating the cue band remains available as a manipulation
check.

Finally, the parameters were verified behaviorally rather than only
acoustically. [PLACEHOLDER — self-test, cells 7-9: own unmodified HRTF,
binaural, fixed-spectrum vs variable-spectrum blocks; report elevation gain for
each and the difference. With intact ears and an unmodified transfer function,
a drop in elevation gain can only reflect the envelope encroaching on the cue,
so an absent or small drop is the direct evidence that the envelope is
transparent to localization.]

% --- supplement ------------------------------------------------------------
% The envelope is a discrete cosine series in dB on a 128-point log-frequency
% grid, coefficient k corresponding to k/10 ripples per octave; the constant
% term and all coefficients at or above 0.5 ripples/octave are zero, a zero
% coefficient contributing 0 dB, i.e. unity gain. Burst-to-burst correlation of
% the detrended log spectra within a trial was r = 0.95, versus r = -0.01 for
% the unmodified stimulus.

% --- what this manipulation can and cannot decide ---------------------------
% Rules out: a mapping tied to the overall spectrum of the training stimulus
%   ("this timbre = that elevation"), since gross coloration now varies trial to
%   trial by as much as elevation itself varies it. Note the strength of this
%   claim depends on the depth: at 3 dB the source only MATCHES the cue at coarse
%   scales, it does not swamp it, so the timbre route is made unreliable rather
%   than uninformative. Phrase it that way in the Discussion; "uninformative"
%   would need the deeper envelope, which M&M's depth threshold rules out.
% Does NOT rule out: a lookup restricted to the preserved 0.5-2 rip/oct band.
%   The pinna notch is a shape-within-band feature, so a remapping of that
%   structure onto space and a learned association between that structure and a
%   response both survive the manipulation. Arguably the two are not
%   distinguishable acoustically at all -- adaptation at the spectral-to-spatial
%   mapping stage IS the learning of new pattern-location associations.
% What would separate perceptual from explicit/procedural learning:
%   - response latency and dual-task cost (explicit strategies are slow and
%     interference-prone)
%   - generalization to untrained locations WITHIN the trained ear: a graded map
%     should interpolate, a lookup should not (cf. condition B)
%   - graded vs catastrophic degradation under band-limiting of the cue region
%   - aftereffects / coexistence on return to the native HRTF
%   - subjective report: does it sound like it comes from there, or is the
%     response reasoned to
% On whether the listener CAN discount the source spectrum here — two reasons
% they can, so a failure is informative rather than a ceiling of the paradigm:
%  (1) Presentation is not strictly monaural. The other ear receives the same
%      source convolved with a cepstrally smoothed (n_keep=4) DTF, so the
%      interaural log-difference cancels the source spectrum exactly while
%      retaining almost all of the trained ear's fine structure. The reference
%      channel is in fact smoother than in natural binaural listening, where
%      both ears carry notches. Its level falls with lateral angle (up to ~9 dB
%      down at 35 deg), so the reference degrades toward the periphery.
%  (2) Source-spectrum discounting is anyway achievable monaurally: models based
%      on the first/second spectral derivative recover direction for sources
%      with locally flat or locally constant-slope spectra (Zakarauskas &
%      Cynader 1993), and current sagittal-plane models use monaural spectral
%      gradients (Baumgartner et al. 2014, already cited in the Intro).
%
% Framing for the Discussion: Macpherson & Middlebrooks show the NORMAL system
% already discounts broad-scale spectral variation. The open question is whether
% a NEWLY acquired spectral-to-spatial mapping inherits that invariance, or is
% initially stimulus-specific and only becomes source-invariant with further
% exposure. That, not "sensory vs cognitive", is what the ripple blocks test.

% --- open items ------------------------------------------------------------
% - Citations to add once a .bib exists:
%     Macpherson & Middlebrooks (2003) JASA 114(1):430-445,
%       doi:10.1121/1.1582174  -- ripple-spectrum probe, the 0.5-2 rip/oct window
%     Zakarauskas & Cynader (1993) JASA 94(3):1323-1331 -- spectral-derivative
%       model, monaural source-spectrum invariance
%     Middlebrooks (1992), Baumgartner et al. (2014) -- already cited in Intro
% - DONE: depth is now reported peak-to-trough as well as rms, and the depth
%   bound is argued explicitly ("Choice of envelope depth").
% - Cue depth varies by ear and subject; the band ratios are scaled from FS's
%   modified left ear at az -20. Decide whether to report a group range instead
%   -- with a shallow envelope the cue-band margin is comfortable everywhere, so
%   a group minimum would be the more convincing number to give.
% - The self-test paragraph is a placeholder until cells 7-9 have been run.
%   If the fixed vs variable difference in elevation gain is not small, the
%   depth argument above has to be rewritten around whatever value survives.
% - FS's two "USO" blocks were run with the old generator (one fixed base
%   texture per session, 0.8 dB variation). They are NOT a variable-spectrum
%   test and should be reported as such, or dropped.


% ---------------------------------------------------------------------------
% DRAFT SEGMENT — the cue manipulation and how the donor is chosen
% Standalone; not \input anywhere yet. Stitch into Methods later.
% Code: hrtf/modify/donor_detail.py       (the composite)
%       hrtf/processing/midline.py        (applied to the measured arc)
%       hrtf/analysis/donor_selection.py  (vanopstal_dtfs / isim / composite_isim)
%
% Code and text now agree: shortlist()/select_donor() gate on ridge_slope <= 0.5
% and rank by |r_match - TARGET_R_MATCH|, both from pair_metrics() on the detail
% domain; DONOR_POOL holds the behaviourally qualified donors and
% MIN_DONOR_ELEVATION_GAIN is the criterion. isim()/vanopstal_dtfs() remain in
% the module as diagnostics and must NOT be reported as similarity.
% ---------------------------------------------------------------------------

\subsection{Spectral cue manipulation}

Participants were trained and tested with a modified version of their own
head-related transfer functions in which the coarse spectral envelope was their
own and the fine spectral detail was another listener's. For each direction and
ear, the log-magnitude spectrum was separated by cepstral truncation into an
envelope, retaining the first four coefficients, and a residual containing
everything finer. The modified spectrum was the participant's own envelope plus
a donor's residual. Own phase was retained and each direction was rescaled to
its original broadband energy, so interaural time differences and broadband
interaural level differences were unchanged and only spectral shape finer than
the envelope scale was altered.

Retaining four coefficients follows Kulkarni and Colburn
% \citet{kulkarni_role_1998}
, who smoothed transfer functions by cepstral truncation and found listeners
could no longer discriminate smoothed from original once roughly sixteen
coefficients were retained; four is well inside the range where the change is
audible and leaves the envelope too coarse to carry an elevation-dependent
pattern. Splitting the spectrum this way rather than substituting the donor's
transfer function whole was verified to decorrelate the spectral-to-spatial map
as thoroughly as a whole substitution, whereas the complementary
composite---donor envelope with own detail---left the native map intact. Detail
therefore carries the elevation cue and the envelope carries timbre, allowing
the cue to be replaced while level, head shadow and individual coloration are
retained.

The manipulation was applied to the measured median-plane transfer functions
and the full directional set was regenerated from them through the
spherical-head model used to build the unmodified set (Sec.~\ref{sec:hrir}), so
that native and modified sets share their interaural time differences exactly
and the envelope is fitted to the pinna response rather than to the pinna
response multiplied by the head shadow.

\subsection{Donor qualification}

Donor transfer functions were taken from a pool of previously recorded
participants. A recording qualified as a donor only if that participant had
themselves localized accurately with their own unmodified transfer function in
the same virtual environment, quantified as the elevation gain of their first
completed localization block (criterion: gain $\geq 0.8$; targets within
$\pm30^\circ$ elevation). The manipulation transplants one listener's spectral
detail to another, and this criterion establishes that the pattern being
transplanted is one a human auditory system demonstrably reads: it is evidence
that the cue works, obtained from the only instrument that can supply it.

The criterion is behavioral because the acoustic alternatives are not
predictive. Among the recordings for which both were available, the depth and
elevation dependence of a donor's own spectral detail---the natural acoustic
index of how much cue it carries---was unrelated to how well that participant
localized with it (Spearman $\rho = +0.18$, $n = 24$); the recording with the
most pronounced fine structure of the entire set belonged to a participant whose
elevation gain was $0.05$. Qualification is applied as an inclusion criterion
only. Accurate localization shows that a pattern is usable; inaccurate
localization is ambiguous between a weak cue, a rendering failure and an
inattentive listener, so an unqualified recording is treated as unproven rather
than as unsuitable.

\subsection{Donor selection}

Among qualified donors, the donor assigned to a participant was the one whose
composite came closest to an intended perturbation size while remaining
unabsorbable by that participant's existing spectral-to-spatial map. Both
quantities were computed between the participant's own and the modified spectral
detail on the median-plane arc, in the 5657--11314\,Hz band, averaged over ears.

Perturbation size was quantified as $r_\mathrm{match}$, the mean over elevations
of the correlation between the participant's own detail at that elevation and
the modified detail at the same elevation. This is the quantity that the
mechanism proposed by Van Wanrooij and Van Opstal
% \citet{vanwanrooij_relearning_2005}
---a correlation between the sensory input and a stored spectral representation
of the listener's own ears---predicts should govern performance. It was targeted
mid-range rather than minimized, because the perturbation is bounded from both
sides: in their monaural experiment, listeners whose cues were not sufficiently
perturbed never reached stable performance, while Trapeau et al.
% \citet{trapeau_fast_2016}
found that larger perturbations produced a greater initial loss and slower, less
complete recovery. The target was set to the median $r_\mathrm{match}$ between
qualified donors, so that the perturbation corresponds to a typical difference
between two people whose own cue is known to work.

Absorbability was quantified as the slope of the regression of the
best-matching own elevation on true elevation. A slope near unity means the
existing map can still read the modified set as a coherent elevation, displaced
rather than replaced, and can accommodate the manipulation as a constant
response bias; this is the configuration that Van Wanrooij and Van Opstal's
non-adapting listeners occupied. Candidates with a slope above $0.5$ were
excluded, and among the remainder the donor closest to the target
$r_\mathrm{match}$ was selected.

Both measures are required. Across all donor--listener pairings available here
they were nearly independent (Spearman $\rho = +0.14$), and of the pairings
whose $r_\mathrm{match}$ indicated a strong perturbation ($\leq 0.3$), $31\%$
nonetheless had a slope above the exclusion threshold---a large change in the
spectral cue that the existing map could still absorb.

Both were computed on the spectral detail rather than on whole transfer
functions. Correlating transfer functions requires the component common to all
directions to be removed by a reference derived from directions other than those
being compared, which for the spectral variability index
% \citep{trapeau_fast_2016}
is a diffuse-field average that the present recordings do not sample. A
reference formed from the median-plane arc itself cannot substitute, since it is
the mean of the same spectra that are then correlated: the residuals sum to zero
by construction, the mean pairwise correlation is fixed at $-1/(N-1)$, and the
resulting index is a constant (measured: $1.02 \pm 0.02$ across 31 recordings,
against $1.056$ predicted for $N = 19$ elevations). The cepstral residual is
subject to no such constraint, because the envelope removed from each direction
is a smoothed copy of that same direction and directions remain independent.

% --- open items ------------------------------------------------------------
% CITATIONS — verified 2026-08-18 unless marked:
%   Van Wanrooij & Van Opstal (2005) J Neurosci 25(22):5413-5424,
%     doi:10.1523/JNEUROSCI.0850-05.2005. "Relearning sound localization with a
%     new ear" (SINGULAR -- the MONAURAL perturbation study). Do NOT confuse
%     with Hofman, Van Riswick & Van Opstal (1998) Nat Neurosci, "with new ears"
%     (plural), the BINAURAL mold study. I_sim definition = their Materials and
%     Methods, "DTF similarity index"; DTF definition = their Eq. 1; the 0.5 and
%     0.3 anchors = their Fig. 7A and 7C; r^2 = 0.94 = their Fig. 6 (n=6).
%     All quoted from the full text, verified.
%   Kulkarni & Colburn (1998) Nature -- cepstral smoothing, discrimination
%     threshold near 16 coefficients. [volume/pages NOT VERIFIED]
%   Trapeau et al. (2016) JASA -- VSI. Cited only if the paragraph declining to
%     use it names it. [volume/pages NOT VERIFIED]
%
% - QUALIFIED POOL = 7 (CO, AGV, SW, VD, AH, FS, FD), from 27 recordings with an
%   own-HRTF block. Every one of the 12 listeners resolves to an admissible
%   donor: 6 in band, 6 widened, no fallbacks. Target r_match 0.58 = median over
%   the 21 qualified-donor pairs (min 0.32, Q1 0.51, Q3 0.66, max 0.77).
%   RECOMPUTE THE TARGET WHENEVER THE POOL CHANGES -- it is defined from the pool.
% - 0.8 falls in the gap between FD at 0.82 and AS at 0.63, so it is not making a
%   fine distinction. State it as a cut in a gap, not as a graded criterion.
% - GS IS A PROBLEM CASE: the only admissible qualified donor for GS is AGV at
%   r_match 0.90, i.e. almost no perturbation. Either GS needs a donor that is
%   not yet qualified, or GS cannot be run under this rule. Check before
%   scheduling GS.
% - Five candidates still have no own-HRTF block (GLK, IM, MSc, NW, UG) -- they
%   have other runs, just none with their unmodified HRTF. One ~60-trial block
%   each would qualify them and widen the pool.
% - CORRECTED: an earlier version of this file called the EG distribution
%   bimodal with a suspicious cluster at zero. That came from
%   vsi_rmse.find_participants() missing 6 subjects (it searched only the flat
%   results/pilot/<id>.pkl and not the nested results/pilot/<id>/<id>.pkl, and
%   preferred the pilot pair where an id existed in both trees, which shadowed
%   the active AS and PF). Both are fixed; with the fuller set the spread is
%   closer to continuous. Do not repeat the bimodality claim.
% - WHY NO I_SIM. Van Opstal's index was tested and rejected for this role: it is
%   the SD down each column of the correlation matrix, and SD is invariant to
%   permutation of the rows, so it cannot distinguish a listener's own ear from
%   their own ear relabelled by 17 deg (measured: identical value 0.542 for
%   shifts of 0, 4, 8 and 17 deg, while the correlation at matching elevation
%   fell 1.00 -> 0.08). It also scores a cue-REMOVED ear (own envelope only,
%   0.395) as less perturbed than a real donor composite (0.272), so it does not
%   separate replacement from removal. Keep this in the file; it is the answer if
%   a reviewer asks why the obvious measure was not used.
% - Detail strength has no repeatability estimate (every arc is measured once),
%   so do not report it to more precision than the ranking requires, and do not
%   claim a difference between adjacent donors.
% - NO REPEATABILITY ESTIMATE exists for any of these measures: every arc is
%   measured once. Re-recording one subject would fix this cheaply, and two
%   recordings HAVE been remade (AS 2026-08-18, GS 2026-07-28) -- but as
%   replacements, not repeats, so they do not serve.
% - Sec.~\ref{sec:hrir} is a placeholder label -- point it at whichever section
%   describes the HRIR recording and the spherical-head azimuth expansion.
% - Regenerate if anything changes: the own-ear and between-listener I_sim
%   values and their n, and the target once set.


\subsection{Statistical analysis}
% TODO — polar error as the primary metric, elevation gain / RMSE / residual
% SD secondary. Transfer measured as within-ear pre-post CHANGE SCORES, not
% the post-only A-vs-D contrast (the mirrored filter carries a foreign-base
% disadvantage that cancels in the change score). Report simple effects of Ear
% within laterality rather than interpreting the Ear x Location interaction.

% --- acoustic methods ------------------------------------------------------

% ---------------------------------------------------------------------------
% DRAFT SEGMENT — HRTF measurement and the virtual HRTF set
% Standalone; not \input anywhere yet. Stitch into Methods later.
% Two \subsection{} blocks. Placed LATE in Methods, where the earmold paper put
% "Binaural recordings" and "Directional transfer functions".
% Code: hrtf/record/recordings.py::record_dome        (acquisition, head tilts)
%       hrtf/record/record_hrir.py                    (orchestration)
%       hrtf/record/processing.py                     (deconvolve, equalize,
%                                                      lowfreq, expand_azimuths)
%       hrtf/record/fit_head_radius.py                (acoustic radius fit)
%       hrtf/record/calibration/calibrate_headphones.py
% Parameters below are read off the code and the recorded params.txt, not the
% module-level interactive defaults (which differ -- head_radius in particular).
% \label{sec:hrir} resolves the placeholder reference in
% methods_donor_selection.tex.
% ---------------------------------------------------------------------------

\subsection{HRTF measurement}\label{sec:hrir}

Head-related impulse responses were measured at the blocked entrance of each
ear canal with MEMS microphones (AMM-3738-T, PUI Devices, Dayton, USA) mounted
in anthropometric earplugs of the PIRATE design
% \citep{denk_pirate_2019}
, fitted to each participant from a size series
% \citep{bau_adapting_2024}
, and supplied with phantom power (PP2B, Millenium). Logarithmic sweeps of 200\,ms duration (120\,Hz to 22\,kHz, 85\,dB SPL) were
presented ten times from each loudspeaker of the midline column and averaged
after alignment. Seven loudspeakers at $12.5^\circ$ spacing sample elevation
too coarsely to resolve the spectral cue, so the arc was recorded three times,
the participant being cued to tilt the head by one third of the loudspeaker
spacing between passes. Interleaving the physical positions in this way yields
19 elevations from $-37.5^\circ$ to $37.5^\circ$ in steps of $4.17^\circ$. A reference recording was then made with the microphones mounted
on a stand at the position the head had occupied and with no listener present,
using identical settings; it captures the response of loudspeakers,
microphones and room in situ. Emitted sweeps were pre-filtered per loudspeaker
with inverse filters obtained from a flat probe microphone, which equalizes
each loudspeaker in level and magnitude response rather than merely matching
them to one another.

Impulse responses were recovered by regularized inversion of the sweep and
each was divided by the reference response of the same loudspeaker. The
quotient is the head-related transfer function (HRTF) in the usual sense: the
pressure at the blocked ear canal entrance relative to the free-field pressure
at the position of the head with the listener absent. These are HRTFs and not
directional transfer functions: no direction-independent component has been
removed. Our earlier work
% \citep{friedrich_adaptation_2025}
reported directional transfer functions, obtained by dividing each transfer
function by the average over all measured transfer functions. That step is
omitted here because only the midline is measured, so an average over the
measured set would be an average over the very directions that are
subsequently compared, and would remove part of the elevation dependence it is
meant to preserve. Each quotient was onset-aligned, windowed with a 2.5\,ms
Hann window and truncated to 512 samples.

\subsection{Virtual HRTF set}

Because only the midline arc was measured, the remaining directions were
synthesized with a rigid-sphere model of the head
% \citep{duda_range_1998}
. The radius of the sphere was fitted for each participant to interaural time
differences measured from a horizontal row of loudspeakers, taken from the
slope of the interaural phase between 200 and 1500\,Hz. The measured arc was then replicated
across azimuth in 25 steps of $4.17^\circ$ from $-50^\circ$ to $50^\circ$,
giving 475 directions, each carrying the spherical-head magnitude shading and
interaural phase difference relative to the frontal direction at the same
elevation. Below 800\,Hz, where the measurement is limited by the sweep
bandwidth and the room, magnitude was replaced by spherical-head values.

One consequence of this construction should be stated plainly. Spectral detail
in the delivered set is measured on the midline and is identical at every
azimuth; azimuth is carried by broadband interaural level and time differences
alone, and the rendering is acoustically faithful to the measurement only at
$0^\circ$ azimuth. Three considerations make this acceptable here. The
elevation-dependent spectral pattern, which is what the manipulation alters
and what the analysis measures, is measured rather than modeled. Every
condition in the design is rendered through the identical expansion, so no
comparison rests on the synthesized component. And because a magnitude-only
edit cannot move a modeled interaural time difference, the native and modified
sets share their interaural cues exactly, which would not be guaranteed had
both been measured. The design accordingly contrasts hemifields rather than
azimuthal acuity, and no claim is made about localization in azimuth.

Two residual corrections were applied. Division by the reference leaves a
small interaural difference at frontal directions, where none should exist;
this was removed by equating broadband energy between the ears over 200\,Hz to
16\,kHz, one scalar per direction, which on the manikin left a frontal
interaural level difference of $0.000 \pm 0.003$\,dB. Headphone equalization
was measured for each participant from three re-seatings of the headphones,
averaged, and inverted as a regularized minimum-phase filter of 256 taps.

% --- open items ------------------------------------------------------------
% - CITATIONS:
%     friedrich_adaptation_2025 -- VERIFIED: Friedrich, P., and Schoenwiesner,
%       M. (2025). J. Acoust. Soc. Am. 157(3), 1543-1553,
%       doi:10.1121/10.0036056.
%     denk_pirate_2019 -- VERIFIED (Zenodo record 2574395). NOTE THE AUTHOR
%       ORDER: Denk, F., Brinkmann, F., Stirnemann, S., and Kollmeier, B.
%       (2019). "The PIRATE: an anthropometric earPlug with exchangeable
%       microphones for Individual Reliable Acquisition of Transfer functions
%       at the Ear canal entrance," Fortschritte der Akustik -- DAGA, Rostock,
%       doi:10.5281/zenodo.2574395. Denk is first author, not Brinkmann.
%     bau_adapting_2024 -- VERIFIED: Bau, D., Moscher, O., and Poerschmann, C.
%       (2024). "Adapting the PIRATE Design for Cost-Effective In-Ear
%       Microphones," Fortschritte der Akustik -- DAGA, Hannover, pp.
%       1604-1606. This is the paper the present hardware follows; it tests
%       AMM-3738-T (PUI) alongside CMM-2718AT (CUI) and SPH1642-HT5H (Knowles)
%       and finds all three near-equivalent, so it supports the choice rather
%       than prescribing it. It is also the source for the XS-XL size series.
%     Duda & Martens (1998) JASA 104(5) -- "Range dependence of the response of
%       a spherical head model". CHECK this is the right one for the model as
%       implemented; if processing.py follows a different formulation, cite
%       that instead. [NOT VERIFIED]
% - MICROPHONE now stated: AMM-3738-T (PUI Devices) + PP2B phantom supply
%   (Millenium is Thomann's house brand -- CHECK whether to name Thomann as the
%   manufacturer instead, and confirm the town for the PUI address line).
%   CONFIRM ALSO that the plugs really are the XS-XL size series rather than
%   individually printed; "fitted to each participant from a size series" is
%   written on the strength of the setup following Bau et al. (2024).
%   The old PUM-3046L-R sentence is gone -- do not let it drift back in.
% - PER-SUBJECT RADII ARE NOT REPORTED. Give the group mean and range once the
%   cohort is closed. Also state what happens on a failed fit (the code falls
%   back to 0.0875 m and stores nothing) -- if any reported participant hit the
%   fallback, that must be said.
% - The head-radius paragraph is now bare by decision (Paul, 2026-08-26): no
%   manikin validation, no Woodworth comparison. If a reviewer asks why the
%   fitted radii (~0.072 m) fall below the Woodworth 0.0875 m, the answer is
%   that the low-frequency phase delay of a real head exceeds its onset delay
%   (Kuhn 1977) and 0.0875 m is the ONSET-delay standard, inappropriate for a
%   phase-based ITD. The fit uses the same estimator the model then imposes.
%   Keep this note; it is the rebuttal, held out of the text on purpose.
% - "limited by the sweep bandwidth and the room" is my reading of why 800 Hz;
%   confirm the actual reason before submission.
% - The 19 elevations come from 7 + 6 + 6 across the three tilt passes (the
%   outermost speaker drops out of the two tilted passes). Not stated in the
%   text -- add if a reviewer queries the count.
% - record_mesm (the newer 7-speaker path) was NOT used for any reported
%   subject. If later participants are recorded on it, this subsection needs a
%   variant, and the two cohorts cannot be pooled without checking.
% - Dome loudspeaker equalization was OFF for subject and reference alike
%   (equalize_dome: False in every params.txt), which is what makes the
%   reference division a clean in-situ calibration. Not stated; probably does
%   not need to be, but it is the answer if asked why no dome EQ is mentioned.


\subsection{Training}
% TODO — cite the training game from Friedrich & Schoenwiesner (2025) (pulse
% rate and inter-pulse delay as feedback, 90 s rounds, scoring, leaderboard)
% and describe only the delta: rendered over headphones through the VAE, one
% ear, and whatever changed in the region sampling.

% ---------------------------------------------------------------------------
% REMOVED 2026-08-26: the second paragraph of the old skeleton, which was still
% about mold 1 / mold 2 and five-day adaptation periods, i.e. the earmold
% paper. Kept here so nothing is lost:
%
%   Participants wore two distinct sets of binaural earmolds consecutively
%   (denoted as mold 1 and mold 2), each over the course of a five-day
%   adaptation period. The acoustical and behavioral impact of the molds were
%   measured, and the trajectories of participants' adaptation were recorded to
%   capture the occurrence of meta-adaptation. To test whether learning a new
%   set of spectral cues interferes with a previously learned mapping, the
%   persistence of adaptation to the first earmolds was measured after
%   participants adapted to the second molds. This measurement was repeated for
%   the second molds after five days without earmolds for comparison.
%
% The previous version of this file is saved as methods.tex.bak_20260826.
% ---------------------------------------------------------------------------
